\documentclass{article}

\usepackage{microtype}
\usepackage{booktabs}
\usepackage{amsmath}
\usepackage{hyperref}
\usepackage{url}
\usepackage{icml2026}

\newcommand{\NA}{\texttt{N/A (not run)}}
\icmltitlerunning{Multi-Agent Peer Review Improvement Loop}

\begin{document}

\twocolumn[
  \icmltitle{Evidence-Grounded Multi-Agent Peer Review:\\
  A Reproducible Author--Reviewer--Chair Improvement Loop}
  \icmltitlerunning{Multi-Agent Peer Review Improvement Loop}

  \begin{icmlauthorlist}
    \icmlauthor{Anonymous Authors}{anon}
  \end{icmlauthorlist}
  \icmlaffiliation{anon}{Anonymous Institution}
  \icmlcorrespondingauthor{Anonymous Authors}{anonymous@example.com}
  \vskip 0.3in
]

\printAffiliationsAndNotice{}

\begin{abstract}
We specify and audit an evidence-grounded, 14-agent paper-review pipeline built
around the repository's nineteen context tasks and strict seven-field review
form. The implementation separates extraction, specialist critique, synthesis,
author rebuttal, and chair output; isolates reviewer and author memories; and
keeps decisions out of paper-time requests. A public 20-case
ReviewBench--ReviewRebuttal seed drives a leakage-safe, versioned calibration
and memory loop. Canonical run \texttt{77771e943444d32dc716} verifies execution,
convergence, best-state restoration, and one sealed test. It uses deterministic
structural signals, not an LLM reviewer, and therefore provides pipeline
evidence rather than reviewer-quality evidence.
\end{abstract}

\section{Scope and Local Contracts}

The system is an offline research instrument for public review histories and
author-authorized feedback. It is not permission to delegate a confidential or
official ICML review to an LLM. ICML 2026 defines reviewer responsibilities and
restricts LLM assistance~\cite{icml2026reviewer,icml2026llm}; the deployment
boundary in Section~\ref{sec:limits} follows that policy.

\paragraph{Nineteen context tasks.}
The authoritative \texttt{prompts.py} map has four groups. \emph{Main Paper}
(4) contains Paper Summary, Supplemental Summary, Introduction, and Related
Works. \emph{Methodology} (6) contains Preliminaries, Framework, Training,
Proofs, Inference and Application, and Method Summary. \emph{Experiments} (7)
contains Datasets, Implementation Details, Evaluation Metrics, Quantitative
Results, Qualitative Results, Ablation Study, and Results Summary.
\emph{Conclusion} (2) contains Limitations and Future Works and Conclusion.
Each extraction request must return separate \texttt{ANSWER} and
\texttt{SOURCES} sections.

The implemented parser wraps each extraction in an immutable
\texttt{Evidence} record with task ID, answer, source strings, provenance, and
a content-derived ID. \texttt{SharedContextStore} appends and deduplicates these
records and returns frozen \texttt{ContextPacket} snapshots. Sources are still
strings: page geometry, typed claims, issues, resolutions, and grounding checks
are not implemented.

\paragraph{Strict final form.}
Only the chair emits \texttt{ReviewOutput}: Soundness, Presentation,
Significance, and Originality are integers 1--4; Overall Recommendation is
1--6; Confidence is 1--5; and Comment is nonempty text. The Markdown parser
requires the exact seven headings, order, blank-line structure, integer types,
and ranges. It does not establish factual accuracy, evidence grounding, or
comment quality. This is a repository-specific subset of the full ICML form.

\section{Team and Information Boundaries}

The registry contains exactly 14 agents: five domain specialists
(CV, Core ML, NLP, RecSys, General), six criterion specialists (Soundness,
Presentation, Significance, Originality, Reproducibility, Ethics), one review
synthesizer, one author, and one chair. A deterministic English/Korean keyword
router selects at most two domain agents, using General only as fallback.

\texttt{ReviewerOrchestrator} runs all nineteen extraction calls in a bounded
pool, restores prompt order, and serializes the resulting packet. Routed domain
and all criterion calls then run in parallel. Their critiques and the synthesis
are free-form strings: prompts request strengths, weaknesses, questions, and
evidence IDs, but no intermediate schema or grounding validator enforces them.
The chair receives serialized context and synthesis and must pass only the
strict final-form parser. With the optional author loop, the author receives
context plus the validated initial review and returns a free-form rebuttal; the
chair then receives that rebuttal and emits a second validated form. Neither
chair call receives a human review or final decision.

The live \texttt{review --state best\_state.json} path uses learned state with
hard role separation. Lexical cue retrieval first excludes memory from the
current forum. Raw training reviews and rebuttals remain internal;
\texttt{memory\_guidance} converts a match into a paper-agnostic checklist.
Domain, criterion, and synthesis requests receive only reviewer guidance, while
the author receives only author guidance. Chair requests alone receive rubric
weights, calibration scale and bias, and state-integrity metadata as advisory
priors. Tests assert self-forum exclusion, raw-text non-disclosure, role
separation, and delivery to both chair calls.

The implementation also prepends an untrusted-document rule: paper and
extracted text cannot override system instructions or reveal unavailable
labels. It does not parse PDFs, detect hidden text, perform OCR comparison,
retrieve semantic evidence, validate free-form critique provenance, or build a
typed issue graph. Those remain required extensions, not reported capabilities.

\section{Public Data and Versioned Update}

OpenReview connects public submissions, reviews, responses, and decisions by
forum~\cite{openreviewnotes,openreviewretrieval}. The canonical smoke does not
claim a fresh API snapshot. \texttt{scripts/build\_real\_seed.py} joins public
paper Markdown from ReviewBench (CC-BY-4.0) with reviews and rebuttal dialogues
from ReviewRebuttal (Apache-2.0)~\cite{reviewbenchdataset,reviewrebuttaldataset}.
It removes author identity fields and keeps cases with a paper, scored review,
rebuttal, and decision. \texttt{data/real/seed\_manifest.json} records source
URLs, source hashes, licenses, selected forum IDs, and output hash. The corpus
SHA-256 begins \texttt{02586f1cfef1}.

The seed has 20 papers, 77 reviews, and 31 rebuttal dialogues: five ICLR 2022,
seven ICLR 2023, and eight NeurIPS 2022, with 15 accept and five reject
decisions. There is no complete RecSys case, and ReviewBench text may be a
revised public version. The implemented SHA-256 forum-ID split is deterministic
and disjoint: 16 train, two development, and two sealed test forums under seed
\texttt{ralphton-icml-real-seed-v1}. It does not yet group PDF/title
near-duplicates or authors. Development and test IDs cannot enter either memory;
the test fingerprint is checked and evaluated once after model selection.

\texttt{LearningState} versions four rubric weights, linear overall-score scale
and bias, prompt-manifest digest, bounded reviewer/author memories, and parent
digest. It neither fine-tunes model weights nor rewrites prompts. On train data,
\texttt{propose\_update} first estimates score calibration. For available final
decisions it computes acceptance probabilities from calibrated overall scores
and applies the clipped residual step
\begin{equation}
 b \leftarrow b + \eta\lambda_{\rm dec}\,
 \operatorname{mean}(y_i-\hat p_i), \qquad \lambda_{\rm dec}=0.25.
\end{equation}
The decision is offline train supervision, never a chair input. A candidate is
accepted only if development field coverage and complete-output coverage do not
regress beyond tolerance and normalized MAE and Brier do not increase beyond
tolerance. Thus decision calibration cannot bypass the per-metric gate.

For version $t$, utility $U_t$ averages available field coverage, complete
coverage, one minus normalized MAE, and one minus Brier score~\cite{brier1950}.
Behavioral delta $D_t$ measures normalized score/probability change and Comment
token-set distance. State delta $R_t$ averages rubric, scale, bias, prompt
digest, and memory change. A plateau satisfies
\begin{equation}
 |U_t-U_{t-1}|\leq\epsilon_U,\quad D_t\leq\epsilon_D,\quad
 R_t\leq\epsilon_R.
\end{equation}
After the minimum iteration, $p$ consecutive plateaus stop training; every
accepted state competes with the best non-regressing state, which is restored
before the single sealed-test call.

\section{Canonical Smoke Result}

Run \texttt{77771e943444d32dc716} identifies eight substantive artifacts in
\path{artifacts/real_seed_v1}. For fixed inputs and configuration their hashes
and run ID are deterministic; \texttt{run\_manifest.json.created\_at} is
execution provenance and changes on rerun. The full-file
\path{dev_review_preview.md} is wrapper-free Markdown that round-trips through
the strict seven-field parser. Prompt-manifest and restored-state SHA-256 values
are, respectively,
\texttt{c0562f49\allowbreak{}640cbe6b\allowbreak{}33e59a6b\allowbreak{}cb2e673a\allowbreak{}2bd7f04c\allowbreak{}b6915a95\allowbreak{}f9805a7c\allowbreak{}f18bcbc2}
and
\texttt{de6d0295\allowbreak{}dbd26723\allowbreak{}e1e6e579\allowbreak{}824362ec\allowbreak{}9dab56f1\allowbreak{}a43eb827\allowbreak{}7c8ffcea\allowbreak{}857330e7}.

The registered configuration used minimum/maximum iterations 3/40, patience 3,
$(\epsilon_U,\epsilon_D,\epsilon_R)=(0.001,0.002,0.002)$, non-regression
tolerance 0.02, learning rate 0.1, decision-calibration weight 0.25, and memory
limit 64. It converged after five iterations and restored version 1 from best
iteration 1, with 64 reviewer and 53 author memory items from train forums.
Four later candidates were rejected; the final three zero-change iterations
satisfied patience.

\begin{table}[t]
\caption{Canonical deterministic smoke. Coverage and complete-output coverage
are 1.0000 throughout; development and test each contain two forums.}
\label{tab:smoke}
\centering
\small
\setlength{\tabcolsep}{4pt}
\begin{tabular}{@{}lrrr@{}}
\toprule
Stage & MAE & Brier & Utility \\
\midrule
Baseline dev & 0.7083 & 0.2863 & 0.8762 \\
Best-state dev & 0.7083 & 0.2796 & 0.8778 \\
Sealed test & 0.3750 & 0.1229 & 0.9426 \\
\bottomrule
\end{tabular}
\end{table}

Development utility rose by 0.001656 solely because Brier improved; MAE was
unchanged. The sealed test has no pre-update counterpart by design. With two
development and two test forums, neither difference supports confidence
intervals, domain comparison, or generalization. No LLM review, extraction
quality study, human annotation, agent ablation, cost study, or robustness
benchmark was run; those results remain \NA{}.

\section{Policy, Limitations, and Reproducibility}
\label{sec:limits}

The system must not ingest confidential submissions or private reviewer
discussion, and its output must not be submitted as an official ICML review.
ICML's LLM and peer-review ethics policies require venue authorization,
confidentiality, integrity, and human accountability
~\cite{icml2026llm,icml2026ethics}. The current no-auth API client and public
seed reduce access risk but do not establish blanket reuse rights. Historical
decisions encode capacity, fashion, disagreement, and social bias; they are
weak calibration supervision, not research truth.

The current evidence core lacks source spans and issue resolution; specialist,
synthesis, and rebuttal outputs are unvalidated strings; no PDF security or
near-duplicate document pipeline exists; routing is lexical; RecSys is absent;
and hosted models may have pretraining contamination or version drift. Human
comment evaluation is costly and subjective. Any quality claim therefore needs
a larger lawful corpus, document-level deduplication, registered LLM baselines
and ablations, blinded domain experts, agreement statistics, robustness tests,
cost reporting, and manifest-linked metric inputs. Formatting follows the
official ICML 2026 author style~\cite{icml2026author}; it does not imply
submission, acceptance, or endorsement.

\paragraph{Conclusion.}
The repository now provides a testable separation of paper-time evidence,
role-specific guidance, offline supervision, calibrated update, convergence,
and sealed evaluation. The canonical artifact establishes that these plumbing
and integrity contracts execute; reviewer competence remains unmeasured.

\clearpage
\bibliographystyle{icml2026}
\bibliography{references}

\end{document}
